\documentclass[a4paper,11pt]{article}
\usepackage[utf8]{inputenc}
\usepackage[T1]{fontenc}
\usepackage[spanish]{babel}
\usepackage{geometry}
\usepackage{array}
\usepackage{colortbl}
\usepackage{xcolor}
\usepackage{booktabs}
\usepackage{enumitem}
\usepackage{fancyhdr}
\usepackage{tikz}

\geometry{left=2cm, right=2cm, top=2.5cm, bottom=2cm}

\definecolor{violeta}{HTML}{7B2D8E}
\definecolor{azulg}{HTML}{3498DB}
\definecolor{rojog}{HTML}{E74C3C}
\definecolor{gris}{HTML}{F2F2F2}
\definecolor{verde}{HTML}{27AE60}
\definecolor{naranja}{HTML}{E67E22}

\pagestyle{fancy}
\fancyhf{}
\fancyhead[L]{\small\textbf{Nuevo Compañeros 1 — U03: La Familia}}
\fancyhead[R]{\small\textbf{Píldora formativa 3.6}}
\fancyfoot[C]{\thepage}

\setlength{\parindent}{0pt}

\begin{document}

\begin{center}
{\LARGE\textbf{PÍLDORA FORMATIVA 3.6}}\\[0.3cm]
{\Large Función comunicativa: decir la hora}\\[0.2cm]
{\normalsize Sección Comunicación (p.39) — Proyectable — 5 diapositivas}\\[0.2cm]
{\small Versión enriquecida secuencial para Fases C1--C4}
\end{center}

\vspace{0.5cm}

\noindent\fbox{\parbox{\dimexpr\linewidth-2\fboxsep-2\fboxrule}{%
\textbf{Contenido para el profesor}\\[0.2cm]
El sistema de la hora en español tiene dos ejes: (1) \textbf{Es / Son}: \textit{Es la una} (singular) vs. \textit{Son las dos, tres...} (plural). (2) \textbf{Y / Menos}: la mitad derecha del reloj usa \textit{y} (y cinco, y diez, y cuarto, y veinte, y veinticinco, y media); la mitad izquierda usa \textit{menos} (menos veinticinco, menos veinte, menos cuarto, menos diez, menos cinco).\\[0.2cm]
Preguntas: \textit{¿Qué hora es?} (hora actual) / \textit{¿A qué hora es + evento?} (hora de un evento). Respuesta con \textit{A la una...} o \textit{A las dos...}\\[0.2cm]
El formato de 12 horas es el estándar en español oral. El de 24 horas se usa en horarios escritos pero se lee en formato de 12: 14:20 = \textit{Son las dos y veinte.}
}}

\vspace{0.8cm}

%% ============================================================
%% DIAPOSITIVA 1
%% ============================================================

\noindent\textcolor{violeta}{\rule{\linewidth}{2.5pt}}

\section*{DIAPOSITIVA 1 — ¿QUÉ HORA ES?}

{\small Fase C1--C2 (Contextualización + Comprensión del modelo) — Técnica: Input flooding situacional con reloj visual}\\
{\small\textit{Principio:} El estudiante ve un reloj analógico con las 12 posiciones etiquetadas. El reloj está coloreado: mitad derecha en verde (Y), mitad izquierda en rojo (MENOS). El color señaliza la regla sin necesidad de explicarla verbalmente.}

\vspace{0.5cm}

\subsection*{En pantalla}

Reloj analógico grande ocupando el 60\% de la pantalla.

\vspace{0.3cm}

\begin{center}
\renewcommand{\arraystretch}{1.4}
\begin{tabular}{|>{\centering\arraybackslash}p{6.5cm}|>{\centering\arraybackslash}p{6.5cm}|}
\hline
\rowcolor{verde!20}
\textbf{\textcolor{verde}{Mitad derecha — Y}} & \cellcolor{rojog!20}\textbf{\textcolor{rojog}{Mitad izquierda — MENOS}} \\
\hline
12 = en punto & 7 = menos veinticinco \\
1 = y cinco & 8 = menos veinte \\
2 = y diez & 9 = menos cuarto \\
3 = y cuarto & 10 = menos diez \\
4 = y veinte & 11 = menos cinco \\
5 = y veinticinco & \\
6 = y media & \\
\hline
\end{tabular}
\end{center}

\vspace{0.3cm}

Debajo del reloj:\\
\textcolor{verde}{\textbf{Zona VERDE = Y}} (sumamos minutos a la hora)\\
\textcolor{rojog}{\textbf{Zona ROJA = MENOS}} (restamos minutos de la hora siguiente)

\vspace{0.5cm}

\subsection*{Instrucciones para el profesor}

\begin{enumerate}[leftmargin=1.5cm, itemsep=4pt]
\item Señale el reloj. Diga: \textit{Este es el reloj de la hora en español.}
\item Señale la mitad verde: \textit{Esta mitad es Y. Cuando el minutero está aquí, decimos y cinco, y diez, y cuarto...}
\item Señale la mitad roja: \textit{Esta mitad es MENOS. Cuando el minutero está aquí, decimos menos veinticinco, menos veinte, menos cuarto...}
\item Diga tres horas señalando el reloj: \textit{Son las tres en punto.} (señale 12) \textit{Son las tres y cuarto.} (señale 3) \textit{Son las cuatro menos cuarto.} (señale 9)
\item Pida que repitan cada hora después de usted (repetición coral).
\end{enumerate}

\newpage

%% ============================================================
%% DIAPOSITIVA 2
%% ============================================================

\noindent\textcolor{violeta}{\rule{\linewidth}{2.5pt}}

\section*{DIAPOSITIVA 2 — ES LA UNA / SON LAS...}

{\small Fase C2 (Comprensión del modelo — contraste es/son) — Técnica: Noticing guiado con contraste visual}\\
{\small\textit{Principio:} El estudiante nota la diferencia es/son. Dos relojes lado a lado: uno marca la 1:00 (ES LA UNA) y otro las 3:00 (SON LAS TRES). El contraste visual es inmediato.}

\vspace{0.5cm}

\subsection*{En pantalla}

Dos relojes lado a lado:

\begin{center}
\begin{tikzpicture}
\node[draw, rounded corners, fill=azulg!10, minimum width=5cm, minimum height=3cm, text width=4.5cm, align=center] at (0,0) {Reloj A: 1:00\\[0.3cm]{\LARGE\textbf{\textcolor{azulg}{ES LA UNA}}}};
\node[draw, rounded corners, fill=naranja!10, minimum width=5cm, minimum height=3cm, text width=4.5cm, align=center] at (7,0) {Reloj B: 3:00\\[0.3cm]{\LARGE\textbf{\textcolor{naranja}{SON LAS TRES}}}};
\end{tikzpicture}
\end{center}

\vspace{0.3cm}

Pregunta: \textit{¿Cuándo usamos ES y cuándo usamos SON?}

\vspace{0.3cm}

Tras clic, aparece la regla:
\begin{itemize}[leftmargin=2cm, itemsep=3pt]
\item \textbf{ES} la una $\rightarrow$ ¡SOLO para la 1:00! (una hora = singular)
\item \textbf{SON} las dos, las tres, las cuatro... $\rightarrow$ Para TODAS las demás (plural)
\end{itemize}

\vspace{0.3cm}

4 relojes pequeños. El estudiante dice si es ES o SON:
\begin{itemize}[leftmargin=2cm, itemsep=3pt]
\item 2:30 $\rightarrow$ \textit{SON las dos y media.}
\item 1:15 $\rightarrow$ \textit{ES la una y cuarto.}
\item 5:00 $\rightarrow$ \textit{SON las cinco en punto.}
\item 1:45 $\rightarrow$ \textit{SON las dos menos cuarto.} (En la zona \textit{menos} se nombra la hora SIGUIENTE.)
\end{itemize}

\vspace{0.5cm}

\subsection*{Instrucciones para el profesor}

\begin{enumerate}[leftmargin=1.5cm, itemsep=4pt]
\item Señale los dos relojes. Lea: \textit{Es la una. Son las tres.}
\item Pregunte: \textit{¿Cuándo usamos ES?} Recoja: para la una, solo una hora. \textit{¿Y SON?} Para las demás.
\item Confirme con la regla que aparece en pantalla.
\item Para los 4 relojes pequeños, señale cada uno y pregunte: \textit{¿ES o SON?}
\item En el último reloj (1:45), si dicen \textit{Es la una menos cuarto}, corrija: \textit{Cuando decimos MENOS, nombramos la hora que VIENE. 1:45 = dos menos cuarto, porque faltan 15 minutos para las dos.}
\end{enumerate}

\newpage

%% ============================================================
%% DIAPOSITIVA 3
%% ============================================================

\noindent\textcolor{violeta}{\rule{\linewidth}{2.5pt}}

\section*{DIAPOSITIVA 3 — SENTENCE BUILDER: CONSTRUID LA HORA}

{\small Fase C2--C3 (Sistematización + Práctica guiada) — Técnica: Sentence Builder (tabla de chunks combinables)}\\
{\small\textit{Principio:} El Sentence Builder muestra las columnas sustituibles del sistema de la hora. El estudiante combina chunks de cada columna para producir expresiones correctas. El weaning off se implementa ocultando columnas progresivamente.}

\vspace{0.5cm}

\subsection*{En pantalla — Sentence Builder}

\begin{center}
\renewcommand{\arraystretch}{1.3}
\begin{tabular}{|>{\centering\arraybackslash}p{2cm}|>{\centering\arraybackslash}p{3cm}|>{\centering\arraybackslash}p{2.5cm}|>{\centering\arraybackslash}p{3cm}|}
\hline
\rowcolor{violeta!15}
\textbf{COL. 1} & \textbf{COL. 2} & \textbf{COL. 3} & \textbf{COL. 4} \\
\hline
Es & la una & & en punto \\
\hline
Son & las dos & \cellcolor{verde!10} y & cinco \\
\cline{4-4}
 & las tres & & diez \\
\cline{4-4}
 & las cuatro & & cuarto \\
\cline{4-4}
 & las cinco & & veinte \\
\cline{4-4}
 & las seis & & veinticinco \\
\cline{4-4}
 & las siete & & media \\
\cline{3-4}
 & las ocho & \cellcolor{rojog!10} menos & veinticinco \\
\cline{4-4}
 & las nueve & & veinte \\
\cline{4-4}
 & las diez & & cuarto \\
\cline{4-4}
 & las once & & diez \\
\cline{4-4}
 & las doce & & cinco \\
\hline
\end{tabular}
\end{center}

\vspace{0.5cm}

\subsection*{Instrucciones para el profesor}

\begin{enumerate}[leftmargin=1.5cm, itemsep=4pt]
\item Señale el Sentence Builder. Diga: \textit{Construid la hora combinando una pieza de cada columna.}
\item Modele: señale Es + la una + y + cuarto: \textit{Es la una y cuarto.}
\item Señale: Son + las cinco + menos + diez: \textit{Son las cinco menos diez.}
\item Pida que los estudiantes construyan 3--4 horas señalando las columnas.
\item \textbf{Weaning off:} Oculte la columna 4 (fracciones). Luego oculte la columna 3 (y/menos). Finalmente, muestre solo un reloj con agujas y los estudiantes dicen la hora completa sin el Builder.
\end{enumerate}

\newpage

%% ============================================================
%% DIAPOSITIVA 4
%% ============================================================

\noindent\textcolor{violeta}{\rule{\linewidth}{2.5pt}}

\section*{DIAPOSITIVA 4 — ¿A QUÉ HORA ES...?}

{\small Fase C3--C4 (Extensión funcional) — Técnica: Extensión funcional con nuevo patrón de pregunta}\\
{\small\textit{Principio:} El estudiante ya sabe decir la hora. Ahora se introduce el patrón para eventos: ¿A qué hora es + evento? con respuesta A las + hora. La cartelera de cine es el contexto motivador.}

\vspace{0.5cm}

\subsection*{En pantalla}

\begin{center}
{\Large\textbf{PARA PREGUNTAR POR UN EVENTO:}}
\end{center}

\vspace{0.3cm}

Patrón:
\begin{itemize}[leftmargin=2cm, itemsep=3pt]
\item Pregunta: \textbf{¿A qué hora es} + [nombre del evento / película]\textbf{?}
\item Respuesta: \textbf{A la una} + fracción / \textbf{A las} + número + fracción.
\end{itemize}

\vspace{0.3cm}

Ejemplo con cartelera:
\begin{itemize}[leftmargin=2cm, itemsep=3pt]
\item \textit{¿A qué hora es Robot Dreams?} $\rightarrow$ \textit{A las cuatro y media.}
\item \textit{¿A qué hora es Padre no hay más que uno 5?} $\rightarrow$ \textit{A las cinco y veinte.}
\end{itemize}

\vspace{0.3cm}

Contraste visual:

\begin{center}
\renewcommand{\arraystretch}{1.5}
\begin{tabular}{|>{\raggedright\arraybackslash}p{6.5cm}|>{\raggedright\arraybackslash}p{6.5cm}|}
\hline
\rowcolor{azulg!10}
\textbf{¿Qué hora es?} & \textbf{¿A qué hora es Robot Dreams?} \\
\hline
\textbf{Son} las cuatro y media. (hora actual) & \textbf{A} las cuatro y media. (hora de un evento) \\
\hline
\end{tabular}
\end{center}

\vspace{0.5cm}

\subsection*{Instrucciones para el profesor}

\begin{enumerate}[leftmargin=1.5cm, itemsep=4pt]
\item Lea el patrón en voz alta. Señale la diferencia: \textit{Cuando decimos la hora actual, empezamos con Son las o Es la. Cuando hablamos de un evento, empezamos con A las o A la.}
\item Modele con el ejemplo de la cartelera.
\item Pida que 2--3 estudiantes formulen preguntas usando la cartelera del libro.
\end{enumerate}

\newpage

%% ============================================================
%% DIAPOSITIVA 5
%% ============================================================

\noindent\textcolor{violeta}{\rule{\linewidth}{2.5pt}}

\section*{DIAPOSITIVA 5 — RELOJ RÁPIDO}

{\small Fase C4 (Producción autónoma) — Técnica: Producción oral coral con ritmo rápido + weaning off}\\
{\small\textit{Principio:} Producción rápida sin tiempo para sobrepensar. El profesor muestra una hora digital en pantalla; los estudiantes la dicen en voz alta lo más rápido posible.}

\vspace{0.5cm}

\subsection*{En pantalla}

Secuencia rápida de 8 relojes digitales (aparecen uno a uno con clic):

\vspace{0.3cm}

\begin{center}
\renewcommand{\arraystretch}{1.6}
\begin{tabular}{|>{\centering\arraybackslash}p{3cm}|>{\raggedright\arraybackslash}p{8cm}|}
\hline
\rowcolor{violeta!15}
\textbf{Hora digital} & \textbf{Respuesta} \\
\hline
\textbf{3:00} & Son las tres en punto. \\
\hline
\textbf{5:15} & Son las cinco y cuarto. \\
\hline
\textbf{1:30} & Es la una y media. \\
\hline
\textbf{8:20} & Son las ocho y veinte. \\
\hline
\textbf{10:45} & Son las once menos cuarto. \\
\hline
\textbf{6:50} & Son las siete menos diez. \\
\hline
\textbf{12:25} & Son las doce y veinticinco. \\
\hline
\textbf{1:55} & Son las dos menos cinco. \\
\hline
\end{tabular}
\end{center}

\vspace{0.5cm}

\subsection*{Instrucciones para el profesor}

\begin{enumerate}[leftmargin=1.5cm, itemsep=4pt]
\item Ronda 1 (con Builder visible): muestre cada hora digital. Los estudiantes dicen la hora en voz alta todos juntos (coral). Confirme.
\item Ronda 2 (sin Builder, solo reloj de Caja 2): repita las mismas horas en orden diferente.
\item Ronda 3 (sin ningún apoyo): señale su reloj de muñeca o diga una hora digital. Los estudiantes responden individualmente (mano alzada).
\end{enumerate}

\end{document}
